%************************************************
\chapter{Variabili}
%************************************************

Prima di iniziare con gli esercizi più semplici, ecco alcuni consigli metodologici per comprendere meglio che uso si sta facendo delle variabili:

\begin{aenumerate}
	\item dare alle variabili nomi brevi, chiari e che aiutino a capirne lo scopo;
	\item assegnare alle variabili solo il loro valore, senza unità di misura o altri dati aggiuntivi;
	\item è possibile aggiungere a fianco delle variabili alcune note per specificare a cosa serve quella specifica variabile e l'eventuale unità di misura del dato in essa contenuto.
\end{aenumerate}

\section{Primi passi con le variabili}

\begin{exercise}
	Realizzare un algoritmo con una sola variabile che inizialmente contiene il numero 3 e, successivamente, il suo valore viene incrementato di 5.
\end{exercise}

\begin{exercise}
	Realizzare un algoritmo con una sola variabile che inizialmente contiene il numero 5, in un secondo passaggio questo valore viene quadruplicato e infine diminuito di 3.
\end{exercise}

\begin{exercise}
	Realizzare un algoritmo con una sola variabile che inizialmente contiene il numero 4, successivamente il suo valore viene incrementati di due per tre volte e, infine, ridotto di 5.
\end{exercise}

\section{Qualche esercizio più serio}

Bla bla bla...

\section{Piccole sfide}

\begin{exercise}
	Bla bla bla bla...
\end{exercise}